%%%%%%%%%%%%%%%%%%%%%%%%%%%%%%%%%%%%%%%%%
% Medium Length Professional CV
% LaTeX Template
% Version 2.0 (8/5/13)
%
% This template has been downloaded from:
% http://www.LaTeXTemplates.com
%
% Original author:
% Trey Hunner (http://www.treyhunner.com/)
%
% Important note:
% This template requires the resume.cls file to be in the same directory as the
% .tex file. The resume.cls file provides the resume style used for structuring the
% document.
%
%%%%%%%%%%%%%%%%%%%%%%%%%%%%%%%%%%%%%%%%%

%----------------------------------------------------------------------------------------
%	PACKAGES AND OTHER DOCUMENT CONFIGURATIONS
%----------------------------------------------------------------------------------------

\documentclass{resume} % Use the custom resume.cls style
\usepackage[left=0.75in,top=0.5in,right=0.75in,bottom=0.5in]{geometry} % Document margins
\usepackage[colorlinks = true,
            linkcolor = blue,
            urlcolor  = blue,
            citecolor = green,
            anchorcolor = blue]{hyperref}

%https://tex.stackexchange.com/questions/392703/change-one-footnote-symbol-from-number-to-asterisk
\usepackage{footnote}
\usepackage{footmisc}
\renewcommand{\thefootnote}{\alph{footnote}}
\newcommand{\astfootnote}[1]{%
\let\oldthefootnote=\thefootnote%
\setcounter{footnote}{0}%
\renewcommand{\thefootnote}{\fnsymbol{footnote}}%
\footnote{#1}%
\let\thefootnote=\oldthefootnote%
}




\newcommand{\tab}[1]{\hspace{.2667\textwidth}\rlap{#1}}
\newcommand{\itab}[1]{\hspace{0em}\rlap{#1}}
\name{Siddhartha Gairola} % Your name
%\address{Microsoft Research Lab, 9, Vigyan 1st floor, Lavelle Road, Bengaluru, Karnataka 560001} % Your address
%\address{123 Pleasant Lane \\ City, State 12345} % Your secondary addess (optional)
\address{\small\textit{E-mail}: \href{mailto:siddhartha.gairola18@gmail.com}{siddhartha.gairola18@gmail.com}}% Your phone number and email
\address{\small\textit{Website:} \href{https://sidgairo18.github.io/}{https://sidgairo18.github.io}}
%\begin{small}\textit{E-mail: }siddhartha.gairola18@gmail.com\end{small}

\pagenumbering{alph}

\begin{document}

%----------------------------------------------------------------------------------------
%	EDUCATION SECTION
%----------------------------------------------------------------------------------------
%\vspace{-2pt}
\begin{rSection}{Education}

{\bf International Institute of Information Technology, Hyderabad} \hfill {Hyderabad, India} 
\\ \textit{Dual Degree BTech and MS by Research in CSE} \hfill {2014 - 2019}
\\ \textit{CGPA: \textbf{8.98 /10}}

{\bf St. Joseph's Academy} \hfill {Dehradun, India} 
\\ \textit{Senior Secondary, ISC; Average: \textbf{93.6\%}} \hfill {2012 - 2013}
\\ \textit{Secondary, ICSE; Average: \textbf{95.28\%}} \hfill {2010 - 2011}

\end{rSection}


%----------------------------------------------------------------------------------------
%	WORK EXPERIENCE SECTION
%----------------------------------------------------------------------------------------

%\vspace{-2pt}
\begin{rSection}{Work Experience}


\begin{rSubsection}{\textbf{Microsoft Research, India} - \textit{Research Fellow}}{Aug 2020 - Present}{}{}
\item Working on problems in healthcare, computer vision and machine learning.
\end{rSubsection}

%\vspace{-4pt}
\begin{rSubsection}{\textbf{Microsoft Research, India} - \textit{Research Intern}}{Jan 2020 - Jul 2020}{}{}
\item Worked on building smartphone based diagnostic solutions for eye diseases using mobile computing, computer vision and machine learning.
\end{rSubsection}

%\vspace{-4pt}
\begin{rSubsection}{\textbf{Adobe Systems, India} - \textit{Research Intern}}{Jun 2019 - Jan 2020}{}{}
\item Worked in \textbf{Media and Data Science Research} (MDSR) lab with Balaji K and Mayur Hemani. 
\item Implemented a deep learning model in \textit{Python, PyTorch} to perform \textbf{Few-shot Guided Segmentation} on images, beat the SOTA by 6\%, published a paper and filed a patent for the same.
\end{rSubsection}

%\vspace{-4pt}
\begin{rSubsection}{\textbf{IIIT Hyderabad} - \textit{Research Assistant}}{May 2016 - May 2019}{}{}
\item Research student at the Center for Visual Information and Technology (CVIT)
\item Worked with Prof. P. J. Narayanan on problems in computer vision, machine learning and deep learning.
\end{rSubsection}

%\vspace{-4pt}
\begin{rSubsection}{\textbf{IIIT Hyderabad} - \textit{Teaching Assistant}}{Aug 2016 - May 2019}{}{}
\item Teaching Assistant for graduate and undergraduate courses - \textit{Digital Logic and Processors, Artificial Intelligence, Digital Image Processing, Computer Vision, and Graphics}.
\item The duties involved taking regular tutorials, setting up assignments and conducting evaluations.
\end{rSubsection}

\begin{rSubsection}{\textbf{Google Summer of Code (GSoC)} - \textit{Student Developer}}{May 2018 - Aug 2018}{}{}
\item Was accepted as a developer for the Google Summer of Code Program for a $2^{nd}$ time with Scilab org.
\item Implemented a DEMO in C/C++ and Scilab as a working example for the MEX Library in Scilab.
\end{rSubsection}

\begin{rSubsection}{\textbf{Google Summer of Code (GSoC)} - \textit{Student Developer}}{May 2017 - Aug 2017}{}{}
\item Was accepted as a developer for the Google Summer of Code Program for the organization Scilab.
\item Implemented a C/C++ wrapper for Matlab MEX-API on current API Scilab.
\end{rSubsection}

\end{rSection}

% Patents
%----------------------------------------------------------------------------------------

\begin{rSection}{Patents}
%\vspace{-4pt}
\begin{rSubsection}{\large \underline{\textbf{}} }{}{}{}
 \item \textbf{Siddhartha Gairola}, Mayur Hemani, Ayush Chopra, Jonas Dahl, Balaji K. (2020), Improved Similarity Propagation for One-Shot and Few-Shot Image Segmentation \textbf{(US 16/906,954)}.
\end{rSubsection}
\end{rSection}



% Publications
%----------------------------------------------------------------------------------------
%\vspace{-7pt}
\begin{rSection}{Publications} \itemsep -2pt

\begin{rSubsection}{\large \underline {\textbf{In Submission}} }{}{}{}
\item \textbf{Siddhartha Gairola}, Murtuza Bohra, Nadeem Shaheer, Navya Jayashankar, Pallavi Joshi, Anand Balasubramaniam, Kaushik Murli, Nipun Kwatra, and Mohit Jain. SmartKC: Smartphone-based Corneal Topographer for Keratoconus Detection. \textit{In Proceedings of the ACM on Interactive, Mobile, Wearable and Ubiquitous Technologies (IMWUT), Volume 5, Issue 4, 2021}.
\end{rSubsection}

\begin{rSubsection}{\large \underline{\textbf{Peer-Reviewed Publications}} }{}{}{}
\item \textbf{Siddhartha Gairola}, Francis Tom, Nipun Kwatra, and Mohit Jain. RespireNet: A Deep Neural Network for Accurately Detecting Abnormal Lung Sounds in Limited Data Setting. \textit{43rd Annual International Conference of the IEEE Engineering in Medicine and Biology Society (EMBC), 2021.}

\vspace{4pt}
\item \textbf{Siddhartha Gairola}, Ayush Chopra, Mayur Hemani, Balaji K. SimPropNet:  Improved Similarity Propagation for Few-Shot Image Segmentation. \textit{International Joint Conference on Artificial Intelligence (IJCAI), 2020}.

\vspace{4pt}
\item \textbf{Siddhartha Gairola}, Rajvi Shah, PJ Narayanan. Unsupervised Image Style Embeddings for Retrieval and Recognition Tasks. \textit{IEEE Winter Conference on Applications of Computer Vision (WACV), 2020}.

\vspace{4pt}
\item Vaibhav Kumar*, Dhruv Khattar*, \textbf{Siddhartha Gairola}*, Yash Kumar Lal*, Vasudev Varma. Identifying Clickbait:  A Multi-Strategy Approach Using Neural Networks. \textit{The 41st International ACM SIGIR Conference on Research \& Development in Information Retrieval, 2018}.

\vspace{4pt}
\item Saumya Rawat*, \textbf{Siddhartha Gairola}\astfootnote{equal contribution}, Rajvi Shah, PJ Narayanan. Find Me a Sky: A Data-Driven Method for Color-Consistent Sky Search and Replacement. \textit{International Conference on Multimedia Modeling (MMM), 2018}.
\end{rSubsection}
\end{rSection}

%	Academic
%----------------------------------------------------------------------------------------
%\vspace{-5pt}
\begin{rSection}{Academic Achievements} \itemsep -5pt
\item Awarded the \textbf{Dean's Research Award} for a paper published at an international conference while in undergraduate, 2018.
\item Awarded the \textbf{Dean's List Award} for 6 semesters straight for being in top 10\% of the class in Monsoon 2015, 2016, 2017 and Spring 2016, 2017, 2018.
\item Qualified for the \textbf{ACM ICPC Asia Onsite Regionals} twice (2015, 2016).
%\item Was a member of the teams \textbf{Pendulum} and \textbf{GOAT} which qualified for the \textbf{ACM ICPC Asia Amritapuri Onsite Regionals} in 2015 and 2016 respectively.
\item \textbf{All India Rank 2262} in \textbf{JEE Advanced 2014} (\textbf{98.5\%} percentile among 150,000 students).
\item \textbf{All India Rank 3856} in \textbf{JEE Mains 2014} (\textbf{99.7\%} percentile among 1.2 million students).
\end{rSection}

%	Academic Service
%----------------------------------------------------------------------------------------
%\vspace{-5pt}
\begin{rSection}{Academic Service} \itemsep -5pt
\item \textbf{Reviewer:} International Conference on Learning Representations (ICLR), 2022; India-HCI (IHCI), 2021. 
\end{rSection}


%	Projects
%----------------------------------------------------------------------------------------
%\vspace{-3pt}
\begin{rSection}{Selected Projects} \itemsep -4pt
%\item \textbf{Hostel Portal:} Developed a hostel management portal using python, Web2py framework. \textit{Spring 2015}

%\item \textbf{Linux C-Shell:} A linux like terminal implemented using \textit{C} language using different concepts like fork, signals, exec, piping. \textit{Monsoon 2015}

%\item \textbf{P2P File Sharing:} Developed an application layer file transfer protocol with support for TCP. It included features such as file upload, file download and indexed searching. (Computer Networks course) \textit{Spring 2016}

\item \textbf{Guided Few-Shot Segmentation}: Implemented a deep learning system to perform few-shot image segmentation with guidance propagation using \textit{Python and PyTorch}. \textit{Monsoon 2019}

\item \textbf{Unsupervised Style Learning}: Implemented a deep learning model in \textit{Python and PyTorch} to learn robust style representations without supervision. Achieved state-of-the-art results for this task and was accepted as a paper (\href{https://github.com/sidgairo18/unsupervised-style-learning}{code}). \textit{Spring 2019}

\item \textbf{Sketching with Style}: Implemented the paper Sketching with Style\footnote{J.  Collomosse,  T.  Bui,  M.  Wilber,  C.  Fang,  and  H.  Jin. Sketching with style:  Visual search with sketches and aesthetic context. In 2017 IEEE International Conference on Computer Vision (ICCV), 2017.} (\textit{ICCV 2017}) in \textit{PyTorch and Python} (\href{https://github.com/sidgairo18/Sketching-with-Style-ICCV-2017-PyTorch-}{code}).  \textit{Monsoon 2018}

\item \textbf{Movie Genre Classification}: Implemented a deep learning model to predict the genres of a movie based on its poster, using \textit{Python and TensorFlow}. \textit{Monsoon 2017}

\item \textbf{Wiki-Search Engine:} Implementation of an efficient and scalable search engine on Wikipedia dump in \textit{Python}. Retrieve top relevant documents based on input query. Handled field as well as phrase queries. Documents were re-ranked based on `tf-idf' measure. For fast retrieval used threading and multi-level indexing.
\textit{Monsoon 2017}

\item \textbf{High Dynamic Range Images:} HDR imaging - generating images with a greater range of luminance levels than which can be achieved by taking only a single photograph with a fixed exposure. Using tone mapping, high-boost filtering and bilateral filtering for improvements.
\textit{Monsoon 2016}

\item \textbf{Tic-tac-toe Bot:} Designed an automated agent bot to play the a modified version of tic-tac-toe game. Using a minimax algorithm and alpha-beta pruning. \textit{Spring 2016}

%\item \textbf{2D Game:} A 2-dimensional shooting game developed using \textit{OpenGl C++}.\textit{Spring 2016}

%\item \textbf{3D Game:} A 3-dimensional game developed using \textit{OpenGl C++}(GLFW library).\textit{Spring 2016}

\end{rSection}

%----------------------------------------------------------------------------------------
%	TECHNICAL STRENGTHS SECTION
%----------------------------------------------------------------------------------------
%\vspace{-5pt}
\begin{rSection}{Technical Skills}

\begin{tabular}{ @{} >{\bfseries}l @{\hspace{1ex}} l }
Working Knowledge &   Bash, C/C++, Python, Matlab, Octave, OpenCV, Android-SDK \\
Software \& Tools & GNU/Linux, HTML, CSS, JavaScript, LaTeX, Excel \\
Deep Learning Frameworks & PyTorch, TensorFlow, Keras, Caffe \\
Past Experience & Java, OpenGL, WebGL, Django, Web2py 
\end{tabular}

\end{rSection}


%	Programming and Open-source
%----------------------------------------------------------------------------------------
%\vspace{-3pt}
%\begin{rSection}{Programming and Open-Source} \itemsep -5pt
%\item \textbf{Programming:} Take part actively in competitive programming contests.\newline
%\textbf{Codechef} (handle: sidgairo18)\newline
%\textbf{Codeforces} (handle: sidgairo18)

%\item \textbf{Open-Source:} Contribute to open source actively. Have made contributions to open source organizations like \textbf{Scilab, LibreOffice and CCExtractor}.
%(\textbf{Github profile} : sidgairo18)
%\end{rSection}

%----------------------------------------------------------------------------------------
%\vspace{-3pt}
%\begin{rSection}{Advanced Courses}
%\itab{\textbf{Core Courses}} \tab{}  %\tab{\textbf{Other Courses}}
%\\ \itab{Computer Programming } \tab{}  \tab{Linear Algebra}
%\\ \itab{Data Structures} \tab{}  \tab{Discrete Mathematics} 
%\\ \itab{Algorithms} \tab{}  \tab{Complexity and Advanced Algorithms} 
%\\ \itab{Operating Systems} \tab{} \tab{Computer System \& Organization}
%\\ \itab{Database Systems} \tab{} \tab{Digital Image Processing}
%\\ \itab{Artificial Intelligence} \tab{} \tab{Computer Vision}
%\\ \itab{Machine Learning} \tab{} \tab{Information Security}
%\\ \itab{Advanced Computer Networks} \tab{} \tab{Systems Network Security}
%\\ \itab{Information Retrieval \& Extraction} \tab{} \tab{Natural Language Applications}
% \\ \itab{Process Control (ongoing)} \tab{} \tab{Electrodynamics}

%\end{rSection}

%----------------------------------------------------------------------------------------
%\begin{rSection}{Extra-Cirrucular} \itemsep -5pt
%\item Play the piano and guitar.
%\item Member of the school football team from 2007-2011.
%\item Member of the IIIT Hyderabad college football team from 2014-2018.
%\end{rSection}


%	References
%----------------------------------------------------------------------------------------
%\vspace{-5pt}
\begin{rSection}{References} \itemsep -5pt
\item \textbf{Prof. P. J. Narayanan}, Director, IIIT-Hyderabad, India \href{mailto:pjn@iiit.ac.in}{(email)}
\item \textbf{Dr. Mohit Jain}, Senior Researcher, Microsoft Research, India \href{mailto:mohja@microsoft.com}{(email)}
\item \textbf{Dr. Nipun Kwatra}, Principal Researcher, Microsoft Research, India \href{mailto:nkwatra@microsoft.com}{(email)}
\item \textbf{Balaji Krishamurthy}, Principal Scientist, Adobe Systems, India \href{mailto:kbalaji@adobe.com}{(email)}
\item \textbf{Mayur Hemani}, Senior Research Scientist, Adobe Systems, India \href{mailto:mayur@adobe.com}{(email)}
\end{rSection}

\end{document}
